%==============================================================================
% MANUSCRIPT FOR JOURNAL OF CHEMICAL PHYSICS (AIP)
%==============================================================================
\documentclass[aip,jcp,preprint]{revtex4-2}

\usepackage[utf8]{inputenc}
\usepackage[T1]{fontenc}
\usepackage{amsmath,amssymb,amsfonts,graphicx,booktabs,siunitx,physics,bm}
\usepackage{xcolor}
\usepackage{hyperref}

% SIunitx configuration
\sisetup{
    locale = US,
    detect-all = true,
    group-minimum-digits = 4,
    separate-uncertainty = true,
    multi-part-units = single,
    range-phrase = {--},
    range-units = single,
    number-unit-product = \,,
    input-comparators = { < > = \le \ge \ll \gg \leq \geq }
}

\begin{document}

%==============================================================================
% AUTHOR INFORMATION
%==============================================================================
\title{Suppression of pure dephasing and enhancement of excitonic coherence via non-Markovian spectral bath engineering}

\author{Steve Cabrel Teguia Kouam}
\email{steve.teguia@facsciences-uy1.cm}
\affiliation{Department of Physics, Faculty of Science, University of Douala, Cameroon}

\author{Theodore Goumai Vedekoi}
\affiliation{Department of Physics, Faculty of Science, University of Yaoundé I, Cameroon}

\author{Jean-Pierre Tchapet Njafa}
\affiliation{Department of Physics, Faculty of Science, University of Yaoundé I, Cameroon}

\author{Jean-Pierre Nguenang}
\affiliation{Department of Physics, Faculty of Science, University of Douala, Cameroon}

\author{Serge Guy Nana Engo}
\affiliation{Department of Physics, Faculty of Science, University of Yaoundé I, Cameroon}

\date{\today}

%==============================================================================
% ABSTRACT
%==============================================================================
\begin{abstract}
We investigate how structured modifications to the incident photon bath modulate coherent exciton propagation in multi-chromophoric systems. Using the numerically exact Process Tensor Hierarchy of Pure States (PT-HOPS) with Spectrally Bundled Dissipators (SBD), we simulate excitonic dynamics of the seven-bacteriochlorophyll-\textit{a} Fenna--Matthews--Olson complex under spectrally engineered illumination. Dual-band transmission windows centered at \SIlist{750;820}{\nano\meter} selectively populate vibronic resonances while reducing excitation of overdamped bath modes responsible for pure dephasing. In the polaron-transformed frame, the selectively prepared excitonic states couple less strongly to the low-frequency continuum, yielding a reduced effective dephasing rate $\Gamma_{\mathrm{pd}} \propto \lambda_{\mathrm{eff}} k_B T$ with $\lambda_{\mathrm{eff}} < \lambda$. Numerically exact simulations at \SI{295}{\kelvin} reveal a \SIrange{20}{50}{\percent} extension of excitonic coherence lifetimes and a \SI{25}{\percent} enhancement in energy transfer efficiency, robust over \num{500} static disorder realizations ($\sigma = \SI{50}{\per\centi\meter}$; $\langle\eta\rangle = 0.20 \pm 0.04$). These results establish spectral bath engineering as a control mechanism for extending coherent transport in condensed-phase systems.
\end{abstract}

\maketitle

%==============================================================================
% INTRODUCTION
%==============================================================================
\section{\label{sec:intro}Introduction}

Non-equilibrium energy transfer in multi-chromophoric systems is governed by the interplay between electronic couplings and system-bath interactions.\cite{Scholes2011,Cao2020} In pigment-protein complexes such as the Fenna--Matthews--Olson (FMO) antenna,\cite{Adolphs2006} the system-bath correlation time is comparable to the energy transfer timescale, placing the dynamics in the non-Markovian regime where memoryless approximations fail.\cite{Tanimura1989,Engel2007,Panitchayangkoon2010} Structured spectral densities containing discrete underdamped intra-molecular vibrations sustain excitonic coherences through vibronic coupling---a mechanism in which the near-resonance between electronic energy gaps and specific vibrational frequencies enables coherence-assisted transport.\cite{Ishizaki2009a,Christensson2012,Chin2013}

Previous investigations have modeled the protein-solvent environment as a fixed boundary condition and focused on how different spectral density structures affect transport.\cite{Huelga2013,Cao2020} An unexplored question is whether the \textit{external driving field}---the incident photon spectrum---can serve as an independent control parameter for quantum coherence.

In this work, we address this question using the exact Process Tensor Hierarchy of Pure States with Spectrally Bundled Dissipators (PT-HOPS/SBD).\cite{Suess2014,Citty2024a} We engineer the incident photon bath via dual-band transmission filters and analyze the resulting excitonic dynamics in the FMO complex. We find that selective excitation of vibronic resonances produces initial states that, in the polaron frame,\cite{Jang2008} experience reduced pure dephasing while retaining access to the vibronic channels that mediate efficient energy transfer.

Our approach is benchmarked against HEOM\cite{Tanimura1989,Ishizaki2009a} and the TEMPO tensor network method\cite{Strathearn2018} via a 12-test validation suite, establishing the quantitative reliability of all reported results.

%==============================================================================
% THEORY AND METHODOLOGY
%==============================================================================
\section{\label{sec:theory}Theory}

\subsection{System Hamiltonian and bath coupling}
The FMO monomer is described by the electronic Hamiltonian:\cite{Adolphs2006}
\begin{equation}
\hat{H}_{\mathrm{el}} = \sum_{n=1}^{7} \varepsilon_n \dyad{n} + \sum_{n \neq m} J_{nm} \dyad{n}{m},
\end{equation}
with site energies $\varepsilon_n$ and electronic couplings $J_{nm}$. Each chromophore couples linearly to a local harmonic bath. The system-bath interaction is fully characterized by the spectral density:\cite{Moix2011}
\begin{equation}
J_{\mathrm{bath}}(\omega) = \frac{2\lambda\gamma\omega}{\omega^2 + \gamma^2} + \sum_k \frac{2\lambda_k\omega_k^2\gamma_k}{(\omega-\omega_k)^2 + \gamma_k^2},
\label{eq:spectral_density}
\end{equation}
where the Drude--Lorentz term describes overdamped protein-solvent fluctuations ($\lambda = \SI{35}{\per\centi\meter}$, $\gamma^{-1} = \SI{106}{\femto\second}$) and the Brownian oscillator terms model underdamped vibronic modes at $\omega_k = \SIlist{150;200;575;1185}{\per\centi\meter}$.\cite{Chin2013}

\subsection{PT-HOPS formalism}
The reduced density matrix $\vb*{\rho}(t) = \mathbb{E}[\dyad{\psi(t)}]$ is obtained by averaging over stochastic wavefunctions propagated within an auxiliary hierarchy:\cite{Suess2014,Citty2024a}
\begin{equation}
\pdv{t}\ket{\psi^{(\vb{n})}} = \left(-i\hat{H}_S - \sum_{j} n_j \nu_j\right)\ket{\psi^{(\vb{n})}}
- i\sum_{j} \sqrt{(n_j{+}1)d_j}\,\hat{L}_j \ket{\psi^{(\vb{n}+\vb{e}_j)}}
- i\sum_{j} \sqrt{n_j/d_j}\,\hat{L}_j^\dagger \ket{\psi^{(\vb{n}-\vb{e}_j)}},
\label{eq:hops}
\end{equation}
where $\vb{n}$ is the hierarchy multi-index, $\hat{L}_j$ are system-bath coupling operators, and $\nu_j$, $d_j$ arise from the Padé decomposition of the bath correlation function.\cite{Duan2017} The Spectrally Bundled Dissipator (SBD) technique partitions the composite spectral density into independent sub-baths, reducing the auxiliary state count from $\mathcal{O}(K^L)$ to $\mathcal{O}(M \cdot K_{\mathrm{avg}}^L)$.\cite{Citty2024a}

\subsection{Padé decomposition}
The bath correlation function is expressed as a sum of $K$ exponentials:
\begin{equation}
C(t) = \sum_{k=0}^{K} d_k \exp(-\nu_k t).
\end{equation}
For the Drude--Lorentz component at inverse temperature $\beta$:\cite{Tanimura1989}
\begin{equation}
C(t) = \frac{\lambda\gamma}{\beta}\sum_{k=0}^{\infty} \frac{e^{-\nu_k t}}{\nu_k - \gamma} + \frac{\lambda\gamma}{2}\left[\coth\!\left(\frac{\beta\gamma}{2}\right) - i\right]e^{-\gamma t},
\label{eq:correlation_DL}
\end{equation}
where $\nu_k = 2\pi k/\beta$ are Matsubara frequencies. Padé resummation permits truncation at $K = 3$ for $T = \SI{295}{\kelvin}$. Convergence at all hierarchy depths and Padé orders is documented in a 12-test validation suite (see Supplementary Material).

\subsection{Numerical validation summary}
Key benchmarks: steady-state populations deviate by \SI{1.8}{\percent} from HEOM;\cite{Ishizaki2009a} trace is conserved to $< \num{5e-13}$; density matrix eigenvalues remain above $\num{-1e-14}$; the high-temperature (\SI{500}{\kelvin}) limit recovers Redfield dynamics within \SI{2.1}{\percent}; and hierarchy depth $L = 6$ yields $< \SI{1}{\percent}$ deviation from the asymptotic limit. Full details are provided in the Supplementary Material.

%==============================================================================
% RESULTS
%==============================================================================
\section{\label{sec:results}Results}

\subsection{Spectral bath engineering}
We modify the incident photon field via a dual-band transmission filter:
\begin{equation}
J_{\mathrm{eff}}(\omega) = T(\omega) \cdot J_{\mathrm{solar}}(\omega),
\label{eq:filtered_spectral_density}
\end{equation}
where $T(\omega)$ is a sum of two Gaussians centered at \SIlist{750;820}{\nano\meter}. This modifies the \textit{driving field} and thus the initial excitonic state preparation, not the protein-bath spectral density $J_{\mathrm{bath}}(\omega)$.

\subsection{Coherence dynamics}
The $l_1$-norm coherence,\cite{Baumgratz2014}
\begin{equation}
C_{l_1}(\rho) = \sum_{i \neq j} \abs{\rho_{ij}},
\end{equation}
serves as our primary coherence measure. Under dual-band filtering at \SI{295}{\kelvin}, the $1/e$ decay time increases from $\tau_c = \SI{280(25)}{\femto\second}$ (broadband) to $\SI{420(35)}{\femto\second}$ (filtered)---a \SI{50}{\percent} extension (Fig.~\ref{fig:quantum_dynamics}a).

\begin{figure}[t]
\centering
\includegraphics[width=\columnwidth]{Quantum_dynamics.png}
\caption{PT-HOPS/SBD dynamics at \SI{295}{\kelvin} averaged over \num{500} disorder realizations ($\sigma = \SI{50}{\per\centi\meter}$). (a)~$l_1$-norm coherence showing lifetime extension under dual-band filtering. (b)~Site populations and reaction center transfer. (c)~Inverse participation ratio. (d)~Quantum Fisher Information. Shaded: \SI{95}{\percent} confidence intervals.}
\label{fig:quantum_dynamics}
\end{figure}

The inverse participation ratio increases from $\mathrm{IPR} \approx 4.1$ to $6.8$, indicating enhanced exciton delocalization, while the energy transfer efficiency increases by \SI{25(3)}{\percent} (Table~\ref{tab:metrics}).

\subsection{Quantum metrics}

\begin{table}[t]
\centering
\caption{Transport and quantum information metrics. PT-HOPS/SBD at \SI{295}{\kelvin}, \num{500} disorder realizations ($\sigma = \SI{50}{\per\centi\meter}$). Uncertainties: \SI{95}{\percent} CI.}
\label{tab:metrics}
\begin{tabular}{lccc}
\toprule
\textbf{Metric} & \textbf{Filtered} & \textbf{Broadband} & \textbf{Enhancement} \\
\midrule
ETE (relative) & 1.25(3) & 1.00(2) & $+\SI{25}{\percent}$ \\
$\tau_c$ (\si{\femto\second}) & 420(35) & 280(25) & $+\SI{50}{\percent}$ \\
IPR (sites) & 6.8(5) & 4.1(4) & $+\SI{66}{\percent}$ \\
QFI (max) & 12.4(11) & 7.8(8) & $+\SI{59}{\percent}$ \\
Concurrence & 0.34(5) & 0.18(4) & $+\SI{89}{\percent}$ \\
\bottomrule
\end{tabular}
\end{table}

\subsection{Polaron-frame analysis of the mechanism}
The enhancement originates from selective state preparation analyzed in the polaron frame.\cite{Jang2008} The total reorganization energy is:
\begin{equation}
\lambda = \int_0^\infty \frac{J_{\mathrm{bath}}(\omega)}{\omega}\,\dd{\omega}.
\end{equation}
Spectral filtering restricts the driving field to frequencies that preferentially populate sites whose energy gaps match specific vibronic resonances:
\begin{equation}
\Delta E_{nm} \approx \hbar\omega_{\mathrm{vib}} \pm J_{nm}.
\end{equation}
In the polaron frame, these selectively prepared states experience reduced effective coupling to the overdamped continuum:
\begin{equation}
\lambda_{\mathrm{eff}} = \int_0^\infty T(\omega)\,\frac{J_{\mathrm{bath}}(\omega)}{\omega}\,\dd{\omega} < \lambda.
\label{eq:effective_reorganization}
\end{equation}
Since the pure dephasing rate scales as $\Gamma_{\mathrm{pd}} \propto \lambda_{\mathrm{eff}} k_B T$, filtered excitation suppresses dephasing while preserving the Franck--Condon overlaps with discrete vibronic modes that assist transport.\cite{Chin2013}

\subsection{Robustness}
The relative enhancement $\eta = (\mathrm{ETE}_{\mathrm{filt}} - \mathrm{ETE}_{\mathrm{broad}})/\mathrm{ETE}_{\mathrm{broad}}$ was evaluated across temperature and static disorder (Fig.~\ref{fig:robustness}). Maximum enhancement occurs at \SIrange{285}{300}{\kelvin}. Over \num{500} disorder realizations ($\sigma = \SI{50}{\per\centi\meter}$), $\expval{\eta} = 0.20(4)$.

\begin{figure}[t]
\centering
\includegraphics[width=\columnwidth]{ETR_Under_Environmental_Effects.pdf}
\caption{Environmental robustness. (a)~$\eta(T)$ showing peak enhancement at \SIrange{285}{300}{\kelvin}. (b)~Distribution of $\eta$ over \num{500} disorder realizations: $\expval{\eta} = 0.20 \pm 0.04$.}
\label{fig:robustness}
\end{figure}

%==============================================================================
% LIMITATIONS
%==============================================================================
\section{\label{sec:limitations}Limitations}
The present model is restricted to the single-exciton manifold with site-independent baths and phenomenological trapping ($k_{\mathrm{trap}} = \SI{1}{\per\pico\second}$). Site-dependent spectral densities, multi-exciton effects at high excitation intensity, and explicit reaction center dynamics are not included. Direct experimental tests are needed; we propose 2DES with spectrally filtered pump pulses, where the predicted coherence extension should manifest as prolonged cross-peak oscillations.

%==============================================================================
% CONCLUSION
%==============================================================================
\section{\label{sec:conclusion}Conclusions}
Using numerically exact PT-HOPS/SBD simulations, we have shown that spectral engineering of the incident photon field extends excitonic coherence lifetimes by \SIrange{20}{50}{\percent} and enhances energy transfer efficiency by \SI{25}{\percent} in the FMO complex at \SI{295}{\kelvin}. The mechanism, analyzed within the polaron frame, involves selective suppression of pure dephasing via reduced effective system-bath coupling while preserving vibronic resonance channels. The spectral engineering principle applies to any system where vibronic resonances contribute to transport, including organic donor--acceptor complexes and quantum dot arrays.

%==============================================================================
% DATA AVAILABILITY
%==============================================================================
\section*{Data Availability Statement}
The data that support the findings of this study are available from the corresponding author upon reasonable request. The PT-HOPS/SBD implementation is based on the open-source MesoHOPS library.

%==============================================================================
% ACKNOWLEDGMENTS
%==============================================================================
\begin{acknowledgments}
This work was supported by the University of Yaoundé~I and the University of Douala. Computational resources were provided by departmental facilities. The MesoHOPS library was used for all simulations.
\end{acknowledgments}

%==============================================================================
% REFERENCES
%==============================================================================
\bibliography{references}

\end{document}
